%% ============================================================================
\section{Application: Guo et al.\ Turbine Engine Data}
\label{sec:application}
%% ============================================================================

To illustrate the sensitivity analysis workflow, we apply it to the masked
failure data from \citet{Huairu-2013}, a $m = 3$ component series system with
$n = 30$ uncensored failures. Each observation records the system failure time
and a candidate set indicating which components could have caused the failure.
This dataset has been used as a benchmark in several masked data
studies~\citep{Huairu-2013, towell2023reliability}.

\paragraph{Standard analysis.}
Fitting the C1-C2-C3 Weibull series MLE yields estimated shapes
$\hat{\mathbf{k}} = (1.26, 1.16, 1.13)$ and scales
$\hat{\vlambda} = (994, 909, 840)$ (the package's canonical fit
\texttt{guo\_weibull\_series\_mle}). The system hazard at the median failure
time $t_0 = 249.5$ is $\hat{h}_T(t_0) = 0.00307$. The large scale estimates
reflect the long lifetimes in this dataset; the shape estimates near 1 suggest
approximately exponential behavior.

\paragraph{Sensitivity sweep.}
We construct a family of $\mathbf{P}$ matrices using the same severity
parameterization as in \Cref{sec:simulations}, with a $3 \times 3$ cyclic
direction matrix and base probability $p_0 = 0.5$. For each severity level,
we refit the Weibull series MLE under the relaxed C2 model with known
$\mathbf{P}(s)$.

\begin{table}[htbp]
\centering
\caption{Sensitivity analysis of the Guo et al.\ dataset ($m = 3$, $n = 30$).
The system hazard $\hat{h}_T(t_0)$ at $t_0 = 249.5$ is stable across all
severity levels; individual scale parameters shift by up to 22\%.}
\label{tab:guo-sensitivity}
\small
\begin{tabular}{c ccc ccc c}
\toprule
$s$ & $\hat{k}_1$ & $\hat{k}_2$ & $\hat{k}_3$ &
      $\hat{\lambda}_1$ & $\hat{\lambda}_2$ & $\hat{\lambda}_3$ &
      $\hat{h}_T(t_0)$ \\
\midrule
0.0 & 1.26 & 1.16 & 1.13 &  994 & 909 & 840 & 0.00307 \\
0.3 & 1.36 & 1.12 & 1.08 &  883 & 923 & 965 & 0.00302 \\
0.5 & 1.34 & 1.10 & 1.15 & 1032 & 918 & 800 & 0.00308 \\
0.8 & 1.35 & 1.08 & 1.16 & 1036 & 933 & 786 & 0.00306 \\
1.0 & 1.36 & 1.07 & 1.16 & 1033 & 945 & 780 & 0.00305 \\
\bottomrule
\end{tabular}
\end{table}

\Cref{tab:guo-sensitivity} confirms the pattern from the simulation study on
real data: the system hazard is remarkably stable (varying by less than 2\%
across the full severity range), while individual component parameters shift
with severity. The scale parameter $\hat\lambda_3$ decreases from 840 to 780
(a 7\% shift) as severity increases, while $\hat\lambda_2$ increases from 909
to 945 (a 4\% shift). Shape parameters show smaller absolute changes but
relative shifts of up to 9\%.

\paragraph{Robustness interval.}
To make the system-level robustness conclusion quantitative rather than
visual, we apply the first-order robustness interval of
\Cref{sec:framework} (function \texttt{ri\_first\_order} in the
package) to the system hazard under a C2 violation. Taking a tolerance of
5\% of the baseline $\hat{h}_T(t_0)$, the estimated local sensitivity index
is near zero, giving a robustness interval of $s \approx 2.07$: the
system-hazard estimate stays within 5\% of its baseline value for C2
severities well beyond the swept range $s \in [0, 1]$. This is consistent
with the near-flat $\hat{h}_T(t_0)$ column of \Cref{tab:guo-sensitivity} and,
because the components are nearly exponential, with the exact all-orders
preservation result for exponential systems. With only $n = 30$ failures the
interval is an illustration of the workflow rather than a precise estimate;
the seed-pinned computation is reproducible via
\texttt{paper/regenerate\_application.R}.

For this dataset, a practitioner can confidently report the system-level
conclusion, that the system hazard at the median failure time is approximately
0.003, without concern about C2 violations. Individual component rankings,
however, are less stable: at $s = 0$, component~3 appears least reliable
($\hat\lambda_3 = 840$, smallest scale), but by $s = 1.0$ component~3 has
$\hat\lambda_3 = 780$ while component~1 has $\hat\lambda_1 = 1033$,
amplifying the gap. Whether this shift matters depends on how the component
estimates will be used.
